\documentclass[11pt]{article}
\usepackage{icmlko}
\begin{document}

\title{줄이는 가중은 학습을 $1.4$ 년 최근으로 민다 --- 거부권이 잡은 것과 못 잡은 것}
\author{월드모델 실험실}
\date{노트 512 · 2026년 8월 3일}
\maketitle

\begin{abstract}
노트 511 이 줄이는 쪽 가중을 노트 446 거부권으로 보류하면서 남겼다 ---
\emph{학습 행 무게가 왜 날짜와 얽히는지는 지금 설명이 없다.} 자료만으로
찾았다.

\textbf{가중이 줄이는 두 도메인이 제일 옛날이다.}

\begin{center}
\begin{tabular}{lrrr}
\hline
도메인 & 학습 행 & 평균 연도 & 무게 \\
\hline
\textbf{만화} & $5{,}647$ & $\mathbf{2013.68}$ & $\mathbf{0.25}$ \\
\textbf{세계애니} & $2{,}648$ & $\mathbf{2015.18}$ & $\mathbf{0.61}$ \\
\hline
도서 & $80$ & $2015.84$ & $1.00$ \\
모바일 & $1{,}559$ & $2018.87$ & $1.00$ \\
웹툰 & $2{,}941$ & $2020.10$ & $1.00$ \\
펀딩 & $2{,}358$ & $2020.72$ & $1.00$ \\
애니 & $2{,}779$ & $2020.88$ & $1.00$ \\
게임 & $259$ & $2021.02$ & $1.00$ \\
\hline
\end{tabular}
\end{center}

줄이는 둘의 평균 연도가 $\mathbf{2014.16}$, 나머지가 $\mathbf{2020.32}$ 다.
그래서 가중을 켜면 \textbf{학습 행 평균 연도가 $2017.566 \to 2018.991$ 로
$+1.425$ 년 밀리고} $2020$ 이전 몫이 $54.4\% \to 46.5\%$ 로 준다.

\textbf{사전 등록 통과.} 노트 446 거부권이 잡은 것은 이것이다 --- 판
$+0.0006$ 중 상당 부분이 \emph{최근성}이고, 날짜 통제 판이 그걸 벗겨
$-0.0005$ 를 남긴다.
\end{abstract}

\begin{figure}[t]
\centering
\includegraphics[width=\textwidth]{figs/f1.pdf}
\caption{\textbf{가} --- 도메인별 학습 연도. \textbf{나} --- 가중이 미는
양. \textbf{다} --- 자 넷의 날짜 처리.}
\end{figure}

\section{그런데 거부권이 못 잡은 것이 있다}

\textbf{집 밖 두 짝의 점수는 이미 날짜를 뺀 값이다.} 이 실험실의 집 밖
채점은 예보와 라벨의 상관에서 \emph{출시일}을 편상관으로 빼고 잰다
(\texttt{part(P, y, day)}). 두 해 동안 그렇게 재 왔다.

그러면 자 넷의 날짜 처리가 이렇게 갈린다.

\begin{center}
\begin{tabular}{lrl}
\hline
자 & 차 & 날짜 \\
\hline
판 & $+0.0006$ & 안 뺌 \\
날짜 통제 판 & $-0.0005$ & 뺌 \\
\textbf{KR 만화} & $\mathbf{+0.0203}$ ($40/40$) & \textbf{뺌} \\
\textbf{비게임 앱} & $\mathbf{+0.0038}$ & \textbf{뺌} \\
\hline
\end{tabular}
\end{center}

\textbf{날짜를 뺀 자 셋 중 둘이 크게 오르고 하나가 조금 내린다.} 즉 이
변경이 사는 것은 최근성만이 아니다 --- 집 밖에서는 날짜를 빼고도
$+0.0203$ 을 산다.

\begin{quote}
\textbf{정직하게} --- 그래도 보류는 유지한다. 노트 446 거부권은 ``판이
오르는데 날짜 통제 판이 내려가면''이라고 썼지 ``집 밖이 좋으면
예외''라고 안 썼다. \textbf{조항을 결과를 보고 고치지 않는다.} 대신
\emph{조항이 겨냥한 것}을 더 정확히 겨냥하는 실험을 설계한다.
\end{quote}

\section{다음 --- 연도 분포를 고정한 가중}

최근성이 섞여 있다면 그것을 빼고 재면 된다. \textbf{가중을 걸되 학습 행의
연도 분포는 원래대로 유지}하는 판본이다 --- 만화 $\cdot$ 세계애니를 줄이되
그만큼 \emph{같은 연도대}에서 채운다(또는 연도별로 재정규화한다).

\begin{quote}
\textbf{사전 등록} --- 연도 분포를 고정한 가중에서 \textbf{판과 날짜 통제
판이 같은 부호}면(둘 다 오르거나 둘 다 $0$ 을 물면) 노트 446 거부권의
전제가 사라지므로, 나머지 자 셋으로 판정한다. 그때도 KR 이 $0$ 밖 양수면
\textbf{채택}.
\end{quote}

이렇게 하면 ``최근성으로 산 것''과 ``구성으로 산 것''이 갈린다. 지금은
섞여 있다.

\section*{남은 것}

\begin{itemize}
\item \textbf{연도 고정 가중} --- 위에 사전 등록.
\item \textbf{씨앗 $80$ 확정}(노트 511) --- 돌고 있다. 날짜 통제 판이
$0$ 을 물면 거부권 전제가 안 서고, 확실한 음수면 이 노트의 기계가 맞다.
\item \textbf{만화가 옛날인 것 자체} --- 학습 $5{,}647$ 행의 평균이
$2013.7$ 이고 $75\%$ 가 $2020$ 이전이다. 노트 505 가 ``만화는 자기 천장의
$2.00$ 배''라고 했는데, 그 천장은 유보($2025+$)에서 잰 것이다.
\textbf{만화는 학습과 채점의 시기가 제일 벌어진 도메인이다.}
\end{itemize}

\end{document}
